\section{Potentials with degenerated ground states, and perturbations thereof}

\begin{definition}[Georgii (19.1)--(19.4): balanced partitions, perturbations, resolving potentials]
    \label{def:dg-19-4}
    \uses{prop:dsb-18-24, def:lo-18-8}
    \lean{MeasureTheory.GibbsMeasure.localGroundEnergyOn, MeasureTheory.GibbsMeasure.energyOutside_eq_localGroundEnergyOn, MeasureTheory.GibbsMeasure.localGroundEnergyOn_iUnion, MeasureTheory.GibbsMeasure.abs_localGroundEnergyOn_sub_le, MeasureTheory.GibbsMeasure.AttainsGroundEnergy, MeasureTheory.GibbsMeasure.attainsGroundEnergy_iff, MeasureTheory.GibbsMeasure.perturb, MeasureTheory.GibbsMeasure.continuous_localGroundEnergyOn_perturb, MeasureTheory.GibbsMeasure.nonnegBoundary, MeasureTheory.GibbsMeasure.nonnegBoundary_eq_frontier, MeasureTheory.GibbsMeasure.groundEnergyGap, MeasureTheory.GibbsMeasure.perturbGap, MeasureTheory.GibbsMeasure.continuous_perturbGap, MeasureTheory.GibbsMeasure.perturbGap_mem_nonnegBoundary, MeasureTheory.GibbsMeasure.perturbGap_zero, MeasureTheory.GibbsMeasure.ResolvesOn, MeasureTheory.GibbsMeasure.Resolves}
    \leanok

    $\lambda$-$\inf_M\Phi_C$ is the essential infimum of the cube interaction on $M$; (19.1) says
    that each member $G_n$ of a stable partition attains the ground energy, taken in Georgii's own
    restatement $\lambda^C(G_n \cap G_\varepsilon(\Phi)) > 0$ for all $\varepsilon > 0$, since over a
    null set the essential infimum is not $+\infty$ in Mathlib. The perturbations (19.2)
    $\Phi^h = \Phi + \sum_{n<N} h_n\Phi^n$, the gap vector (19.3) $g(h) = (\lambda\text{-}\inf_{G_n}
    \Phi^h_C - \lambda\text{-}\inf_G \Phi^h_C)_n$, continuous with values in $\partial\mathbb R^N_+$
    (which is the frontier of the nonnegative orthant) and $g(0) = 0$ under (19.1); Definition
    (19.4): $\Phi^1,\dots,\Phi^{N-1}$ resolve the partition on a neighbourhood $U$ of $0$ if $g$
    maps $U$ onto a neighbourhood of $0$ in $\partial\mathbb R^N_+$ with a continuous left inverse
    (Georgii's homeomorphism), and $\lambda^C(G_\delta(\Phi^h))$ is bounded below on $U$ for each
    $\delta > 0$.
\end{definition}

\begin{proposition}[Georgii Comment (19.5)]
    \label{prop:dg-19-5}
    \uses{def:dg-19-4}
    \lean{MeasureTheory.GibbsMeasure.gapModel, MeasureTheory.GibbsMeasure.gapHomeomorph, MeasureTheory.GibbsMeasure.perturb_eq_of_mem, MeasureTheory.GibbsMeasure.perturbGap_eq_gapModel, MeasureTheory.GibbsMeasure.localGroundEnergyOn_perturb_iUnion_eq_localGroundEnergy, MeasureTheory.GibbsMeasure.localGroundStates_perturb_inter_eq, MeasureTheory.GibbsMeasure.resolvesOn_of_const_on, MeasureTheory.GibbsMeasure.resolves_of_const_on, MeasureTheory.GibbsMeasure.exists_measure_localGroundStates_pos_of_finite}
    \leanok

    (1) Bounded $\Phi^n$ with $\Phi^n_C = c_n + \delta_n$ on $G_n$ and $c_n$ on $G \setminus G_n$,
    $\delta_n \neq 0$, resolve the partition: $g(h) = (h_n\delta_n - \min_k h_k\delta_k)_n$ on the
    neighbourhood $\sum_k |h_k| b_k < \varepsilon/2$, an explicit homeomorphism onto
    $\partial\mathbb R^N_+$, and $G_\delta(\Phi^h) \cap G_n = G_\delta(\Phi) \cap G_n$ where
    $g_n(h) = 0$; disjointness of the partition is used explicitly. (2), first sentence: for a
    finite state space with $\lambda$ charging every point, condition (19.4)(ii) holds for every
    family with the bound $(\min_x\lambda\{x\})^{2^d}$; the metric criterion of its second
    sentence, which Georgii does not prove, is not formalised.
\end{proposition}
\begin{proof}
    \leanok
\end{proof}

\begin{theorem}[Georgii Corollary (19.7) and Example (19.8)]
    \label{thm:dg-19-7}
    \uses{def:dg-19-4, thm:dg-19-6, thm:dsb-18-25, lem:lo-ocean-tail}
    \lean{MeasureTheory.GibbsMeasure.measurableSet_setOf_mem_latticeOceanPart, MeasureTheory.GibbsMeasure.measure_setOf_latticeOceanPart_nonempty_eq_of_map, MeasureTheory.GibbsMeasure.measure_setOf_latticeOceanPart_nonempty_eq_inv, MeasureTheory.GibbsMeasure.exists_gibbsMeasure_refinement, MeasureTheory.GibbsMeasure.blumeCapelCube, MeasureTheory.GibbsMeasure.blumeCapelField, MeasureTheory.GibbsMeasure.blumeCapelEnergy_eq_zero_iff, MeasureTheory.GibbsMeasure.localGroundEnergy_blumeCapel, MeasureTheory.GibbsMeasure.localGroundStates_blumeCapel, MeasureTheory.GibbsMeasure.isPStablePartition_blumeCapel, MeasureTheory.GibbsMeasure.isPStablePartition_blumeCapelFine, MeasureTheory.GibbsMeasure.attainsGroundEnergy_blumeCapel, MeasureTheory.GibbsMeasure.resolves_blumeCapel, MeasureTheory.GibbsMeasure.isRSymmetric_of_const, MeasureTheory.GibbsMeasure.IsRSymmetric.image_cubeRefl, MeasureTheory.GibbsMeasure.blumeCapelCube_cubeRefl, MeasureTheory.GibbsMeasure.blumeCapelField_cubeRefl, MeasureTheory.GibbsMeasure.isCubeConfining_blumeCapelCube, MeasureTheory.GibbsMeasure.isRSymmetric_blumeCapelPartition, MeasureTheory.GibbsMeasure.blumeCapelSpinFlip, MeasureTheory.GibbsMeasure.blumeCapelCube_onCube_blumeCapelSpinFlip, MeasureTheory.GibbsMeasure.blumeCapelField_onCube_blumeCapelSpinFlip, MeasureTheory.GibbsMeasure.isPStablePartition_blumeCapelSplit, MeasureTheory.GibbsMeasure.eventually_exists_gibbsMeasure_family_blumeCapel, MeasureTheory.GibbsMeasure.eventually_three_le_encard_extremePoints_blumeCapel}
    \leanok

    Corollary (19.7), stated for exactly what its proof uses: a Gibbs measure $\mu$ carrying an
    ocean of $\bigcup_\ell G_{n,\ell}$ almost surely, with the sub-partition stable and its members
    transported into each other by symmetries preserving $\mu$, splits by conditioning into $k_n$
    mutually singular Gibbs measures with $\mu(A_{n,\ell}) = 1/k_n$, each carrying an ocean of its
    own pattern and with $\mu_{n,\ell}(0 \in \xi^0(G_{n,\ell})) = \mu(0 \in \xi^0(\bigcup G))$, an
    equality where Georgii has an inequality; it applies to the output of (18.25) and of (19.6).
    Example (19.8), the Blume--Capel model on $\{-1,0,1\}$: the local ground states are the three
    constant configurations, $\{G_1, G_2\}$ with $G_2 = \{(\pm\pm)\}$ and its refinement are stable
    and balanced, and the self-potential $\sum_{i\in C}\sigma_i^2$ resolves $\{G_1, G_2\}$ by (19.5)(1)
    with $c_1 = 4$, $\delta_1 = -4$. Its conclusion
    $|\operatorname{ex}\mathcal G_\Theta(\beta\Phi + \beta h(\beta)\Phi^1)| \geq 3$ is now Theorem
    \ref{thm:dg-19-6} applied to $\{G_1,G_2\}$ and the resolving self-potential, followed by (19.7)
    splitting the second phase along the spin flip $\sigma \mapsto -\sigma$: the model-specific
    inputs are the reflection invariance of $\Phi_C$ and $\Phi^1_C$ (both are a bond sum and a site
    sum over the unit square), the $r$-symmetry of the partitions (their members are constant
    configurations), the boundedness of $\Phi_C$ over the finite state space, and the spin flip
    $0 \leftrightarrow 2$ as a $\lambda$-preserving cube transformation preserving $\Phi_C$ and
    $\Phi^1_C$ and exchanging $G_{2,1}$ with $G_{2,2}$. Corollary (19.7) uses only condition
    (18.23)(a)(ii) of the refinement, whereas \texttt{IsPStablePartition} bundles that with
    (18.23)(a)(i); the refinement $\{G_{2,1}, G_{2,2}\}$ is therefore recorded as a stable partition
    of the local ground states of $\Phi - \Phi^1$, the member $h = -1$ of the family (19.2), whose
    local ground states are exactly $G_2$. Not formalised: Examples (19.9)--(19.10), and Bricmont
    and Slawny's theorem that the unperturbed $\beta\Phi$ has a unique shift-invariant Gibbs
    measure, which Georgii quotes without proof.
\end{theorem}
\begin{proof}
    \leanok
\end{proof}

\begin{theorem}[Sperner's lemma and the Knaster--Kuratowski--Mazurkiewicz lemma, Georgii (19.12) and Appendix 19.A]
    \label{thm:dg-19-12}
    \lean{Sperner.IsCell, Sperner.IsAdmissible, Sperner.IsPanchromatic, Sperner.panchromaticCells, Sperner.odd_card_panchromaticCells, Sperner.exists_isPanchromatic, Sperner.card_doors_mod_two, Sperner.pivot, Sperner.card_panchromaticCells_face, Sperner.even_card_of_involutive_of_fixedPointFree, Sperner.card_mod_two_eq_card_filter_fixed, Sperner.bary, Sperner.bary_mem_stdSimplex, Sperner.abs_bary_vertex_sub_le, nonempty_iInter_of_isClosed_of_stdSimplex_subset_iUnion, exists_mem_stdSimplex_forall_exists}
    \leanok

    Sperner's lemma on Kuhn's triangulation of $\ell\Delta^d$ in staircase coordinates
    ($0 \leq y_0 \leq \dots \leq y_{d-1} \leq \ell$, cells indexed by a lattice point and a
    permutation): an admissible colouring (each vertex coloured by an index of its support) has an
    odd number of panchromatic cells, by the door-to-door parity argument through a pivot
    involution on doors, the fixed doors being the panchromatic cells of the face $\alpha_d = 0$;
    $\ell \geq 1$ is needed. Lemma (19.12): closed $K_1,\dots,K_N \subseteq \mathcal P_N$ covering
    the simplex with $K_n \cap \{\alpha_n = 0\} = \emptyset$ have $\bigcap_n K_n \neq \emptyset$,
    by Georgii's (19.A4) (colour the grid points by covering members, take a limit of panchromatic
    cells of mesh $2/\ell$); and the selection principle in the form (19.13) uses. Erratum:
    Georgii's (19.A1) is stated for the cliques of the graph on the grid in which two points are
    adjacent when they differ by $\pm\ell^{-1}$ in exactly two coordinates, claiming an $n$-clique
    is a simplex of the decomposition; for $N \geq 4$ the $N$-cliques are only the sets
    $\{y + e_m\}$ and $\{z - e_m\}$, and an admissible colouring of $\mathcal P_{4,2}$ with no
    panchromatic $4$-clique exists (checked by hand and by enumeration), while the classical
    statement on the standard triangulation, which is what (19.A4) uses, holds.
\end{theorem}
\begin{proof}
    \leanok
\end{proof}

\begin{theorem}[Georgii \S19.2: Propositions (19.13)--(19.15) and Theorem (19.6)]
    \label{thm:dg-19-6}
    \uses{thm:dg-19-12, def:dg-19-4, prop:dg-19-5, thm:dsb-18-25, thm:lo-18-17, lem:lo-ocean-tail, prop:lo-18-12, thm:exg-4-22}
    \lean{MapClusterPt.fst, MapClusterPt.snd, Filter.Tendsto.mapClusterPt_prod, Potential.gibbsSpecificationOfSigmaFiniteAdmissible_le_of_abs_hamiltonian_sub_le, MeasureTheory.abs_integral_sub_integral_le_of_le_smul, MeasureTheory.GibbsMeasure.cubeInteractingSet, MeasureTheory.GibbsMeasure.abs_hamiltonian_cubePotential_sub_le, MeasureTheory.GibbsMeasure.tendsto_dist_action_cubeGibbsSpec, MeasureTheory.GibbsMeasure.tendsto_dist_action_cubeGibbsSpec_of_eventually, MeasureTheory.GibbsMeasure.exists_mem_stdSimplex_forall_exists_measure_le, MeasureTheory.GibbsMeasure.le_cond_measure_zero_mem_latticeOceanPart, MeasureTheory.GibbsMeasure.exists_gibbsMeasure_family_of_isCountablyCompact, MeasureTheory.GibbsMeasure.le_encard_extremePoints_invariantG_shiftGroup_of_separated, MeasureTheory.GibbsMeasure.le_encard_extremePoints_of_latticeOceanPart, MeasureTheory.GibbsMeasure.isQuasiGibbsianWith_of_isGibbsMeasure_cubeGibbsSpec, MeasureTheory.GibbsMeasure.le_measure_mem_latticeOceanPart_of_forall_le_pow, MeasureTheory.GibbsMeasure.patternWeight_le_of_measure_diff_eq_zero, MeasureTheory.GibbsMeasure.le_crossingBound, MeasureTheory.GibbsMeasure.IsCubeConfining.add_of_abs_le, MeasureTheory.GibbsMeasure.exists_correspondence_of_patternWeight_le, MeasureTheory.GibbsMeasure.exists_gibbsMeasure_family_of_patternWeight_le, MeasureTheory.GibbsMeasure.simplexDisadvantage, MeasureTheory.GibbsMeasure.simplexDisadvantage_mem_nonnegBoundary, MeasureTheory.GibbsMeasure.exists_delta_box_subset_image, MeasureTheory.GibbsMeasure.exists_perturbation_map, MeasureTheory.GibbsMeasure.eventually_forall_patternWeight_perturb_le, MeasureTheory.GibbsMeasure.eventually_exists_gibbsMeasure_family_of_resolvesOn}
    \leanok

    Proposition (19.13), for a family $(\Gamma_\alpha)_{\alpha \in \mathcal P_N}$ of sets of random
    fields with (i) $\{(\alpha,\mu) : \mu \in \Gamma_\alpha\}$ countably compact in the product of the
    Euclidean and the local topology (Mathlib's countable compactness, literally Georgii's ``every
    sequence has a cluster point''), (ii) $\mu(0 \in \xi^0(G,\cdot)) \geq 1 - c$ on $\Gamma_\alpha$
    with $(N+1)c < 1$, and (iii) the boundary condition $N\mu(0 \in V(G_n,\cdot)) < \mu(0 \in V(G,\cdot))$
    when $\alpha_n = 0$: Step 1 is the Knaster--Kuratowski--Mazurkiewicz selection with the ratios
    $\mu(0 \in V(G_n,\cdot))/\mu(0 \in V(G,\cdot))$, the sets $K_n$ being sequentially closed by the
    product cluster-point lemmas; Step 2 conditions on the tail events $A_n$ through (7.7)(b),
    giving $N+1$ mutually singular shift-invariant Gibbs measures with
    $\mu_n(0 \in \xi^0(G_n,\cdot)) \geq 1 - (N+1)c$, and $|\operatorname{ex}\mathcal G_\Theta| \geq
    N+1$ from measures separated by disjoint events. Proposition (19.14): for a family of
    $C$-potentials $\Psi^\alpha$ continuous in $\alpha$ along the simplex, with $t(G,\Psi^\alpha)
    \leq t$ and $(N+1)t(E^C\setminus G_n,\Psi^\alpha) < 1 - t$ on the faces, and a common confining
    sequence, the sets $\Gamma(\alpha)$ of Gibbs measures of $\Psi^\alpha$ meeting the bound satisfy
    (i)--(iii) with $c = 4z(t)$: nonemptiness from $\mathcal G_0(\Psi^\alpha)$ by (18.10), (18.12) and
    (18.22); (i) from (4.19) for a $C$-potential --- proved here through the two-sided kernel
    domination $\gamma^\Psi_\Lambda \leq e^{2|\beta|D}\gamma^\Phi_\Lambda$ when the Hamiltonians differ by
    at most $D$, since the tree's (4.19) covers absolutely summable potentials only and (17.18)(iv)
    allows unbounded interactions --- with (4.12), (4.9), (4.17) and the semicontinuity (18.9)(3);
    (ii) is (18.17) from the chessboard estimate alone, for any Gibbs measure of the cube
    specification. (19.15): the map $\alpha \mapsto h(\alpha)$ with $g(h(\alpha)) = 2\delta(1 -
    (\max_k\alpha_k)^{-1}\alpha_n)_n$ from the left inverse of the resolving map on a box of the
    boundary, with $\delta$ made smaller instead of Georgii's neighbourhood, and $t_\beta \to 0$
    uniformly on the neighbourhood by (18.9)(4). Theorem (19.6): for standard Borel $E$, a stable
    partition into $r$-symmetric $G_n$ satisfying (19.1) and resolving potentials, for $\beta$ large
    some $h(\beta)$ in any prescribed neighbourhood of $0$ makes $\beta\Phi^{h(\beta)}$ carry $N+1$
    mutually singular shift-invariant Gibbs measures with $\mu_n(\xi^0(G_n,\cdot) \neq \emptyset)
    = 1$ and $\mu_n(0 \in \xi^0(G_n,\cdot)) \geq 1 - 4(N+1)z(t)$, so
    $|\operatorname{ex}\mathcal G_\Theta(\beta\Phi^{h(\beta)})| \geq N+1$; Georgii's
    $(N+1)t < 4Nz(t) < 1$ is derived from $4(N+1)z(t) < 1$, and his $h(\beta) \to 0$ is the statement
    quantified over the neighbourhood. Georgii's (19.13)(iii) is stated with the denominator
    cleared, and (19.14)(I) relative to the simplex, which is forced since $\alpha \mapsto h(\alpha)$
    is continuous only along it.
\end{theorem}
\begin{proof}
    \leanok
\end{proof}
